%! Logarithmic Functions — Unit 2, Lesson 2.11 — Lesson Plan
\documentclass[10pt]{article}
\usepackage{precalculus-article}
\usepackage{precalculus-boxes}
\usepackage{graphicx}
\graphicspath{{images/}}
\newcommand{\TallMath}[1]{$\displaystyle #1\rule[-1.4em]{0pt}{3.2em}$}
\newcommand{\CourseName}{Precalculus}
\newcommand{\SchoolYear}{2026--2027}
\newcommand{\MeetingLength}{55 minutes}
\newcommand{\UnitNumberName}{Unit 2: Exponential and Logarithmic Functions \quad}
\newcommand{\LessonNumberName}{Lesson 2.11: Logarithmic Functions}

\begin{document}
\begin{center}
    {\Large\bfseries \CourseName: \SchoolYear \\
    {\small\bfseries \UnitNumberName \LessonNumberName}}
\end{center}

\begin{tcolorbox}[colback=sky, colframe=navy, boxrule=0.9pt, arc=2mm,
    left=2mm, right=2mm, top=1.5mm, bottom=1.5mm]
    \textbf{Primary Objective:} Students identify key characteristics of logarithmic functions
    in general form $f(x) = a\log_b x$, including domain, range, vertical asymptote, end
    behavior, monotonicity, and concavity, and use the additive transformation identification
    test to determine whether a function is logarithmic.
    \hfill\textit{\small Skills: AP 3.A \quad LO 2.11.A}
\end{tcolorbox}

\renewcommand{\arraystretch}{1.6}
\begin{skillbox}[Priority Ideas \& Skills]{goldbox}
\begin{tabularx}{\linewidth}{@{}X X@{}}
\textbf{AP Skills} &
\textbf{Key Understandings} \\
\midrule
\textbf{AP 3.A} --- Describe characteristics of a function with varying levels of precision. &
\begin{itemize}[leftmargin=1.2em,itemsep=2pt,topsep=0pt]
  \item Domain of $f(x) = a\log_b x$ is $x > 0$; range is all reals. (EK 2.11.A.1)
  \item Logarithmic functions are always monotone and always concave up or always concave
        down; no extrema except on a closed interval; no inflection points. (EK 2.11.A.2)
  \item End behavior: $\lim_{x\to 0^+} a\log_b x = \pm\infty$ and
        $\lim_{x\to\infty} a\log_b x = \pm\infty$; vertical asymptote at $x=0$. (EK 2.11.A.4)
  \item If $g(x) = f(x+k)$ has inputs proportional over equal output intervals, then
        $f$ is logarithmic; the shifted version $g$ does NOT have this property. (EK 2.11.A.3)
\end{itemize} \\
\end{tabularx}
\end{skillbox}

\begin{skillbox}[{Vocabulary, Concepts, \& Theorems}]{greenbox}
\renewcommand{\arraystretch}{1.4}
\begin{tabularx}{\linewidth}{@{} >{\bfseries\color{navy}}p{0.30\linewidth} X @{}}
logarithmic function (general form) & $f(x) = a\log_b x$, where $b>0$, $b\ne1$, $a\ne0$;
                                      the inverse family of general-form exponential functions
                                      with initial value 1. \\
domain & All real numbers greater than zero: $(0, \infty)$.
         Cannot take the log of zero or a negative number. \\
range & All real numbers $(-\infty, \infty)$. \\
vertical asymptote & The line $x = 0$ (the $y$-axis) for $f(x) = a\log_b x$;
                     the graph approaches but never touches it. \\
end behavior & How $f(x)$ behaves as $x\to 0^+$ and as $x\to\infty$;
               direction depends on the sign of $a$ and whether $b>1$ or $0<b<1$. \\
monotone increasing / decreasing & Always increasing or always decreasing on its full domain;
                                   logarithmic functions have no local extrema. \\
concave up / concave down & The graph curves in one direction only; logarithms have no
                            inflection point. \\
additive transformation identification test & If $g(x) = f(x+k)$ has inputs proportional over
                            equal-length output intervals, then $f$ is logarithmic. \\
\end{tabularx}
\end{skillbox}

\begin{skillbox}[Activate Prior Knowledge \& Spiral Review]{sky}
    Spiral review (warm-up) covers: evaluating $\log_2 16$, $\log_3 1$, and $\log_5 5$ to
    solidify basic logarithm mechanics; completing a table for $g(x) = 3^x$ and identifying
    which inputs are valid for its inverse $f(x) = \log_3 x$ (domain restriction); and a
    true/false question about whether $f(x) = 2\log_5 x$ has a $y$-intercept (False ---
    $x = 0$ is not in the domain). These connect to EK 2.11.A.1 and EK 2.11.A.4.
\end{skillbox}

\begin{skillbox}[Hook]{sky}
    Write on the board: \textbf{The Richter scale uses $M = \log_{10}(I)$, where $I$ is
    earthquake intensity.}

    Ask: ``A magnitude-1 quake has intensity $I = 10$. What intensity gives magnitude 2?
    Magnitude 3?'' Have students fill a small table. Then: ``Magnitude 2 is how many times
    more intense than magnitude 1?'' (10$\times$.) ``Magnitude 3 vs.\ magnitude 1?'' (100$\times$.)
    Pose: \textbf{``If the output (magnitude) grows by 1 each time the input (intensity)
    multiplies by 10, what type of function is this? What inputs make sense?''} This motivates
    domain $(0,\infty)$, unbounded range, end behavior, and the additive growth of logs.
\end{skillbox}

\begin{skillbox}[Lesson]{sky}
\begin{multicols}{2}
\textbf{Part 1 (10 min): Domain, Range, and Asymptote (EK 2.11.A.1, A.4).}

Start with $f(x) = \log_{10} x$. Ask: ``Can we evaluate $f(0)$?'' (No --- log of zero is
undefined.) ``$f(-5)$?'' (No --- log of a negative is not real.) Conclude: domain is
$x > 0$. For range: show $f(0.001) = -3$, $f(1\,000\,000) = 6$ --- outputs cover all reals.
EK 2.11.A.4: as $x \to 0^+$, $f(x) \to -\infty$; as $x \to \infty$, $f(x) \to +\infty$
(for $b>1$, $a>0$). The $y$-axis is a vertical asymptote because the domain excludes
$x \le 0$.

\columnbreak

\textbf{Part 2 (10 min): Monotonicity and Concavity (EK 2.11.A.2).}

Since $f(x) = \log_b x$ is the inverse of $g(x) = b^x$, it inherits structural properties.
Ask: ``Is $g(x) = 2^x$ always increasing?'' (Yes.) ``So is $f(x) = \log_2 x$ increasing or
decreasing?'' (Always increasing.) For $0 < b < 1$, the exponential is decreasing, so the
log is also always decreasing. Ask: ``Is there any peak or valley?'' (No --- no extrema on
the natural domain, no inflection points.) This is the content of EK 2.11.A.2.

\textbf{Part 3 (10 min): End Behavior and Identification Test (EK 2.11.A.3, A.4).}

Write the two limit statements; walk through cases $b>1$ vs.\ $0<b<1$, $a>0$ vs.\ $a<0$.
Then introduce EK 2.11.A.3: the additive transformation test. Show the notes table:
given $f(x) = \log_3 x$ and the inputs $x = 1,3,9,27,81$ (outputs $0,1,2,3,4$), the
inputs are proportional (multiply by 3) over equal output intervals. Then form $g(x) = f(x+2)$
and check the new inputs --- they are \emph{not} proportional, confirming EK 2.11.A.3.
\end{multicols}
\end{skillbox}

\begin{skillbox}[Active Monitoring]{redbox}
\begin{itemize}[leftmargin=1.2em,itemsep=2pt,topsep=0pt]
  \item After Part 1: cold-call ``What is the domain of $h(x) = -3\log_2 x$?'' Expect $x > 0$.
        Students who say ``all reals'' need domain reinforcement.
  \item Watch for $x \ge 0$ instead of $x > 0$ --- the asymptote at $x = 0$ means $x = 0$
        is excluded. Ask: ``What happens to $\log_2 x$ as $x \to 0^+$?'' (Goes to $-\infty$.)
  \item After Part 2: ask ``Does $f(x) = \log_2 x$ have a maximum on its natural domain?''
        Expect ``No.'' Students who say yes are confusing a closed-interval restriction with
        the general domain. Cite EK 2.11.A.2.
  \item During Part 3: distinguish the sign of $a$ from the value of $b$ when reading limits.
        Write a 2$\times$2 table ($a>0$/$a<0$ vs.\ $b>1$/$0<b<1$) on the board.
\end{itemize}
\end{skillbox}

\begin{skillbox}[Group Work \& Differentiation]{redbox}
\begin{itemize}[leftmargin=1.2em,itemsep=2pt,topsep=0pt]
  \item \textbf{Tier R (Remediate):} Domain, range, asymptote of $f(x) = \log_{10} x$;
        complete a values table ($x = 0.1,\,1,\,10,\,100,\,1000$); identify invalid $x$-values.
  \item \textbf{Tier A (Apply):} For $f(x) = 2\log_3 x$: domain, range, asymptote, end
        behavior limits ($x \to 0^+$ and $x \to \infty$); monotonicity and concavity in words.
  \item \textbf{Tier E (Extension):} Additive transformation identification test using
        $f(x) = \log_3 x$ table and $g(x) = f(x+k)$; apply test to a mystery function table.
\end{itemize}
\end{skillbox}

\begin{skillbox}[Individual Work \& Assessment]{redbox}
Exit ticket (2 questions, 1 page):
\begin{enumerate}[leftmargin=1.5em,itemsep=2pt,topsep=2pt]
  \item State the domain and range of $f(x) = 3\log_4 x$. Write the equation of the vertical
        asymptote. (Tests EK 2.11.A.1, A.4.)
  \item For $f(x) = \log_2 x$: fill in the two end-behavior limits; state whether $f$ is
        increasing or decreasing; state whether $f$ is concave up or concave down.
        (Tests EK 2.11.A.2, A.4.)
\end{enumerate}

Sort by result: students missing \#1 need domain/asymptote work (EK 2.11.A.1); students
missing \#2 need end behavior and concavity (EK 2.11.A.2, A.4).
\end{skillbox}

\begin{skillbox}[Reinforcement \& Extension]{goldbox}
\textbf{Homework (5--6 problems):} domain/range/asymptote for three functions in $a\log_b x$
form; end behavior and monotonicity for $g(x) = -\log_2 x$; additive transformation
identification problem; written explanation of why logarithms have no local extrema;
AP-style MC identifying a non-characteristic of $\log_3 x$; extension comparing end behavior
of $\log_2 x$ vs.\ $\log_{1/2} x$.

\textbf{Preview of Lesson 2.12:} Transformations of logarithmic functions --- horizontal and
vertical shifts, reflections, and how domain and asymptote shift under $f(x) = a\log_b(x-h)+k$.
\end{skillbox}

\end{document}
